\section{An elementary finite upper bound}
\label{sec:upper}

We record a coarse quantitative form of the classical block-and-focus
induction; see van der Waerden~\cite[pp.~142--145]{vanDerWaerden1998}.
The recursion below proves finiteness directly, involves only previously
defined integers, and applies to every positive number of colors.

\begin{theorem}
\label{upper:finite}
Define $F(r,1)=1$ and $F(r,2)=r+1$ for every integer $r\ge1$.
For $k\ge3$, with $r$ and $k$ fixed, define
\begin{equation}
 m_0=1,\qquad
 m_{t+1}=2m_tF(r^{m_t},k-1)\quad(0\le t<r),\qquad
 F(r,k)=m_r.
 \label{upper:recursion}
\end{equation}
Then, for all positive integers $r,k$,
\[
 W_r(k)\le F(r,k)<\infty.
\]
The boundary parameter values are exactly
\[
 W_r(1)=1,\qquad W_r(2)=r+1,\qquad W_1(k)=k.
\]
\end{theorem}

\begin{proof}
The recursion is well founded: construct all $F(r,k)$ at a given
length $k$ after constructing them for every number of colors at
length $k-1$.  Each of the $r$ steps in
Equation~\eqref{upper:recursion} then uses finite positive integers.

A one-term progression is a singleton, so $W_r(1)=1$.
Any $r$-coloring of $[r+1]$ has two positions of the same color, and
those positions form a two-term progression with positive difference.
An injective coloring of $[r]$ shows that $W_r(2)=r+1$.
Also, in one color $[k]$ contains a $k$-term progression and $[k-1]$
does not, proving $W_1(k)=k$.

It remains to prove the asserted upper bound for $k\ge3$.
We induct on $k$, simultaneously for all $r$, with the induction
statement that every $r$-coloring of an interval of $F(r,k)$
consecutive integers contains a monochromatic $k$-term progression.
The preceding paragraph proves the cases $k=1,2$.
Fix $k\ge3$, assume the statement at length $k-1$ for every color
count, and fix $r\ge1$.

For the inner induction, a \emph{focus with $t$ colors} in a colored
interval consists of a point $f$ of that interval and positive integers
$d_1,\ldots,d_t$ such that all the points
\[
 f-sd_i,\qquad 1\le s\le k-1,\quad 1\le i\le t,
\]
belong to the interval, and, for each $i$, the $k-1$ points indexed by
$i$ have a common color $a_i$, with $a_1,\ldots,a_t$ pairwise
distinct.  In an interval containing no monochromatic $k$-term
progression, the focus itself must have a color different from each
$a_i$: otherwise the corresponding points, in increasing order,
together with $f$ form such a progression of difference $d_i>0$.

We claim, by induction on $t\in\{0,\ldots,r\}$, that every
$r$-coloring of an interval of length $m_t$ containing no monochromatic
$k$-term progression has a focus with $t$ colors.  For $t=0$ this just
means choosing the single point in an interval of length $m_0=1$.

Assume the claim for $t<r$, and put
\[
 m=m_t,\qquad B=F(r^m,k-1),\qquad m_{t+1}=2mB.
\]
By translating, we may work with an avoiding coloring $\chi$ of
$[2mB]$.  Partition it into $2B$ consecutive blocks of length $m$,
and associate to block $j$ the word
\[
 w_j=\bigl(\chi((j-1)m+1),\ldots,\chi(jm)\bigr).
\]
There are at most $r^m$ different words.  Apply the outer induction
hypothesis at length $k-1$ and color count $r^m$ to the word coloring
of the first $B$ block indices.  We obtain integers $b,\Delta\ge1$
such that
\begin{equation}
 w_b=w_{b+\Delta}=\cdots=w_{b+(k-2)\Delta},\qquad
 b+(k-2)\Delta\le B.
 \label{upper:matching-blocks}
\end{equation}
Since $k-2\ge1$, we have $\Delta\le B-1$.  Thus the next index
\begin{equation}
 q=b+(k-1)\Delta
 \le B+\Delta\le2B-1<2B
 \label{upper:next-block}
\end{equation}
is still one of the allocated blocks.  Its word need not agree with
the preceding words.

The common word in Equation~\eqref{upper:matching-blocks}, regarded
as a coloring of $[m]$, contains no monochromatic $k$-term progression,
since such a progression would also occur inside block $b$.
Apply the inner induction hypothesis to this word.  Let $f\in[m]$
be its focus, with differences $d_1,\ldots,d_t$ and distinct arm
colors $a_1,\ldots,a_t$.  Let $a_0$ be the color at its position $f$;
as observed above, $a_0$ is distinct from all $a_i$.

Take position $f$ in block $q$ as the new focus, namely
\[
 f'=(q-1)m+f.
\]
For each $i\in\{1,\ldots,t\}$ and $s\in\{1,\ldots,k-1\}$,
\begin{align*}
 f'-s(m\Delta+d_i)
 &=\bigl(b+(k-1-s)\Delta-1\bigr)m+(f-sd_i).
\end{align*}
The block on the right is one of the matching blocks, and the
within-block position $f-sd_i$ belongs to $[m]$.
Consequently these points all have color $a_i$ for fixed $i$.
There is one more arm:
\[
 f'-sm\Delta
 =\bigl(b+(k-1-s)\Delta-1\bigr)m+f,
 \qquad 1\le s\le k-1,
\]
whose points all have color $a_0$.
All $t+1$ arm differences are positive, all arm points are in the
interval, and their colors are pairwise distinct.  Together with
Equation~\eqref{upper:next-block}, this proves that $f'$ is a focus
with $t+1$ colors, completing the inner induction.

At $t=r$, an avoiding coloring of $[m_r]$ would have a focus whose
arms use $r$ distinct colors.  Its own color would have to avoid all
of them, which is impossible in an $r$-coloring.  Hence every such
coloring contains a monochromatic $k$-term progression, completing
the outer induction.  In particular $W_r(k)$ exists and is at most
$F(r,k)$.
\end{proof}
